%% Example CV - Carlos Jadir Ocampo de la Torre
%% Comprehensive overview of all competencies, experience, and achievements.
%% Use as a master reference when tailoring role-specific CVs.
%%
%% SETUP: Run /setup to populate this with your actual information.
%% Compile with: lualatex main_example.tex
%% (pdflatex often fails on modern MiKTeX with fontawesome5 font-expansion errors.)

\documentclass[11pt,a4paper,sans]{moderncv}
\moderncvstyle{banking}
\moderncvcolor{blue}

% Force the name and section headings to render in moderncv blue (color1).
% Default banking leaves them black: moderncvstylebanking.sty's \colorlet
% copies (not aliases) the pre-scheme accent colour, so the name colours are
% frozen before \moderncvcolor runs. Re-let them after. \namefont is the hook
% every name-style macro routes through, so this also works on moderncv 2.3.1
% (Debian/Ubuntu apt), which has no \firstnamestyle/\lastnamestyle at all.
\renewcommand*{\namefont}{\fontsize{34}{36}\bfseries\upshape}
\colorlet{firstnamecolor}{color1}
\colorlet{lastnamecolor}{color1}
\colorlet{namecolor}{color1}
\renewcommand*{\sectionstyle}[1]{{\sectionfont\color{color1}#1}}

\usepackage[utf8]{inputenc}
% pdflatex fallback only (the documented engine is lualatex, which skips this
% branch). Without T1 font encoding pdflatex builds accented letters with
% \accent, and the PDF text layer stores them decomposed - `e` + U+0300 rather
% than U+00E8 - so an ATS keyword match on "Genève" fails while the page looks
% right. moderncv 2.5 loads T1 itself under pdflatex; 2.3.1 (Debian/Ubuntu apt)
% does not. \ifpdftex comes from iftex, which every moderncv version loads.
\ifpdftex\usepackage[T1]{fontenc}\fi
% moderncv loads hyperref itself in an \AtEndPreamble hook, so \hypersetup
% must go in an \AtEndPreamble of our own: on moderncv < 2.4 a top-level
% \usepackage{hyperref} clashes with the class's own
% \RequirePackage[unicode]{hyperref}. From 2.4.0 the class passes its options
% through \PassOptionsToPackage instead, which is what removes that clash.
\AtEndPreamble{\hypersetup{
    colorlinks=true,
    linkcolor=blue,
    filecolor=magenta,
    urlcolor=blue,
    pdftitle={Carlos Jadir Ocampo de la Torre - CV},
    % Keep pdfpagemode=UseNone: this block runs after moderncv's own
    % \AtEndPreamble (moderncv.cls sets pdfpagemode there), so a FullScreen
    % value here would win and open every CV in fullscreen presentation mode.
    pdfpagemode=UseNone,
}}
\usepackage[scale=0.80]{geometry}
\usepackage{import}

% personal data
\name{Carlos Jadir}{Ocampo de la Torre}
% If you have no address to list, DELETE this whole line. \address{}{}{} fails
% with "There's no line here to end" on every moderncv version.
\address{Zihuatanejo, Guerrero, México}{}{}
\phone[mobile]{+52 7443260318}
\email{cjocampo27@outlook.com}
\extrainfo{\href{https://www.linkedin.com/in/carlos-ocampo-510918253}{LinkedIn}, \href{https://github.com/Carlex27}{GitHub}}

\begin{document}
\makecvtitle

\vspace{6pt}
\small{Desarrollador full-stack y móvil con experiencia en APIs, aplicaciones Flutter y automatización de documentos. Ha integrado bases de datos, dispositivos IoT y modelos locales para resolver necesidades operativas.}

\section{Competencias técnicas}
\begin{itemize}
  \item \textbf{Backend y datos:} Python, FastAPI, Java, Spring Boot, Node.js, PostgreSQL, MySQL y Firebase.
  \item \textbf{Web y móvil:} Dart, Flutter, TypeScript, React, Next.js y Vite.
  \item \textbf{Automatización e IoT:} Ollama, procesamiento de PDF, ESP32, NFC, APIs REST y WebSockets.
  \item \textbf{Herramientas:} Git, GitHub, Docker, Postman y Linux/Bash.
\end{itemize}

\section{Experiencia y proyectos}
\begin{itemize}
  \item{\cventry{Oct 2026}{Desarrollador Full-Stack}{Sistema de clasificación arancelaria de acero}{}{}{
  \begin{itemize}
    \item Implementó búsqueda, clasificación y sugerencia de fracciones arancelarias y NICO.
    \item Automatizó la extracción de certificados de molino en PDF e integró modelos locales en Ollama para formatos no conocidos.
    \item Desarrolló el historial de certificados y la exportación a Excel.
  \end{itemize}}}

  \item{\cventry{Ene 2026 - Jun 2026}{Lead Full-Stack \& IoT Developer}{ZIQUA}{}{}{
  \begin{itemize}
    \item Diseñó la integración entre el sensor ESP32, Firebase y la aplicación Flutter para monitorear pH, temperatura y turbidez.
    \item Desarrolló un módulo predictivo para mantenimiento preventivo y funciones de alertas y gestión del servicio.
  \end{itemize}}}

  \item{\cventry{Feb 2025 - Dic 2025}{Líder de Proyecto y Desarrollador Principal}{SALL-E}{}{}{
  \begin{itemize}
    \item Lideró el desarrollo técnico y la estrategia comercial del proyecto, ganador regional de InnovaTecNM 2025.
    \item Desarrolló la aplicación Flutter con Firebase para GPS y diagnóstico de carga, e implementó autenticación vehicular con tarjetas NFC encriptadas.
  \end{itemize}}}
\end{itemize}

\section{Educación}
\begin{itemize}
  \item{\cventry{Ago 2022 - Jun 2027}{Ingeniería en Sistemas Computacionales}{Instituto Tecnológico Nacional de México, Campus Costa Grande}{}{}{En curso. Graduación prevista en junio de 2027.}}
\end{itemize}

\section{Idiomas}
Español nativo. Inglés B2, conversacional y técnico.

\section{Reconocimientos}
Primer lugar regional y finalista nacional de InnovaTecNM 2025. Tercer lugar regional de Hackatec 2026.

\end{document}
